\section{Placement}
\label{sec:placement}

Once the chunks have been defined, the model must decide where each chunk is
placed on the pallet. More precisely, for each chunk it must choose the pallet
row on which the chunk is dropped and the horizontal position at which the chunk
starts. This is represented by two decision variables:

\vspace{0.5cm}
\begin{lstlisting}
array[1..n_chunks] of var 1..boxes_per_row: x_start;
array[1..n_chunks] of var 1..pallet_rows: y_row;
\end{lstlisting}
\vspace{0.5cm}

The variable \texttt{y\_row[n]} denotes the row assigned to chunk \(n\), while
\texttt{x\_start[n]} denotes the column where that chunk begins.

A first requirement is that every chunk must fit within the horizontal limits of
the pallet. Since chunk \(n\) starts at column \texttt{x\_start[n]} and has
length \texttt{chunk\_len[n]}, its last occupied column is
\texttt{x\_start[n] + chunk\_len[n] - 1}. This position must not exceed the
number of columns available on the pallet:

\vspace{0.5cm}
\begin{lstlisting}
constraint forall(n in 1..n_chunks)(
  x_start[n] + chunk_len[n] - 1 <= boxes_per_row
);
\end{lstlisting}
\vspace{0.5cm}

The second requirement is that two chunks must not occupy the same pallet cells.
Each chunk can be seen as a horizontal rectangle of height \(1\): it starts at
\((\texttt{x\_start[n]}, \texttt{y\_row[n]})\), has width
\texttt{chunk\_len[n]}, and has height \(1\). Therefore, the non-overlap
condition can be expressed directly by using 
\begin{lstlisting}
constraint forall(i, j in 1..n_chunks where i < j)(
  y_row[i] != y_row[j]                              \/
  x_start[i] + chunk_len[i] <= x_start[j]          \/
  x_start[j] + chunk_len[j] <= x_start[i]
);
\end{lstlisting}
\vspace{0.5cm}

or with the global constraint \texttt{diffn}:

\vspace{0.5cm}
\begin{lstlisting}
include "globals.mzn";   % for diffn

constraint diffn(x_start, y_row, chunk_len, [1 | n in 1..n_chunks]);
\end{lstlisting}
\vspace{0.5cm}

The constraint \texttt{diffn(x, y, dx, dy)} enforces that the
rectangles defined by origins \((\texttt{x[i]}, \texttt{y[i]})\) and
dimensions \((\texttt{dx[i]}, \texttt{dy[i]})\) do not overlap. In this model, all
rectangles have height \(1\), so \texttt{diffn} ensures that chunks assigned to
the same row occupy disjoint horizontal intervals.

Together with the chunking constraints, these placement constraints already
define a valid pallet layout. The chunking constraints ensure that all
\texttt{n\_boxes} boxes are assigned to chunks, while the placement constraints
ensure that those chunks fit inside the pallet and do not overlap.

However, one useful structural property is still only implicit: each pallet row
must receive exactly \texttt{chunks\_per\_row} chunks. Nevertheless, leaving
this property implicit makes the solver discover it through search.

To make this information explicit, we can add to the model a global cardinality
constraint:

\vspace{0.5cm}
\begin{lstlisting}
constraint global_cardinality(
  y_row,
  [r | r in 1..pallet_rows],
  [chunks_per_row | r in 1..pallet_rows]
);
\end{lstlisting}
\vspace{0.5cm}

The constraint \texttt{global\_cardinality(arr, vals, counts)} enforces that
each value \texttt{vals[i]} appears in \texttt{arr} exactly
\texttt{counts[i]} times. In this case, every pallet row from \(1\) to
\texttt{pallet\_rows} must appear exactly \texttt{chunks\_per\_row} times
in \texttt{y\_row}. As a result, the solver knows from the beginning how many
chunks must be assigned to each row. Once a row has received the required number
of chunks, no further chunks can be placed there, which reduces unnecessary
branching during search.

\section{Release and Plate Collision}
\label{sec:release}

When a chunk is released, the gripper opens one of its two plates. The
opening side sweeps the adjacent pallet row: if that row already contains a
box overlapping the chunk along \(x\), the plate collides with it. A binary
variable encodes the opening side, with \(0\) meaning the plate opens
upward (toward row \(r-1\)) and \(1\) downward (toward row \(r+1\)):

\vspace{0.5cm}
\begin{lstlisting}
array[1..n_chunks] of var 0..1: plate;
\end{lstlisting}

First, it introduces an
auxiliary variable that materialises the target row of each plate:

\vspace{0.5cm}
\begin{lstlisting}
array[1..n_chunks] of var 0..boxes_along_y+1: target;
constraint forall(n in 1..n_chunks)(
  target[n] = y_row[n] - 1 + 2*plate[n]
);
\end{lstlisting}
\vspace{0.5cm}

When \texttt{plate[n] = 0} the expression yields \texttt{y\_row[n] - 1};
when \texttt{plate[n] = 1} it yields \texttt{y\_row[n] + 1}. The two
opening directions collapse into a single arithmetic equality, replacing
the reified disjunction of the initial formulation. The domain is widened
to \([0, \texttt{boxes\_along\_y}+1]\) so that the boundary values \(0\)
and \(\texttt{boxes\_along\_y}+1\) explicitly represent "off the pallet";
they never match a real row index and so they encode a free edge without
any special case.

Second, the temporal precedence is built directly into the quantifier
range. Because the position of a chunk in the belt stream coincides with
its release index, only pairs with \(pred_chunk < current_chunk\) describe a potential
collision (\(pred\_chunk\) is already on the pallet when \(current\_chunk\) is released). The
range \texttt{pred\_chunk in 1..current\_chunk-1} is a compile-time parameter, so the solver
generates only the relevant pairs:

\vspace{0.5cm}
\begin{lstlisting}
constraint forall(current_chunk in 2..n_chunks, pred_chunk in 1..current_chunk-1)(
  not(
       y_row[pred_chunk] = target[current_chunk]
    /\ x_start[pred_chunk] < x_start[current_chunk]  + chunk_len[current_chunk]
    /\ x_start[current_chunk]  < x_start[pred_chunk] + chunk_len[pred_chunk]
  )
);
\end{lstlisting}
\vspace{0.5cm}

The body becomes a direct negation of a single conjunction: chunk \(pred\_chunk\)
cannot share the target row of \(current\_chunk\) while overlapping it along \(x\). No
reification on the plate value (encoded by \texttt{target}), no
reification on the temporal ordering (encoded by the quantifier range).
A single propagator per pair, where the previous formulation needed
several.  